% Options for packages loaded elsewhere
% Options for packages loaded elsewhere
\PassOptionsToPackage{unicode}{hyperref}
\PassOptionsToPackage{hyphens}{url}
\PassOptionsToPackage{dvipsnames,svgnames,x11names}{xcolor}
%
\documentclass[
  letterpaper,
  DIV=11,
  numbers=noendperiod]{scrartcl}
\usepackage{xcolor}
\usepackage{amsmath,amssymb}
\setcounter{secnumdepth}{-\maxdimen} % remove section numbering
\usepackage{iftex}
\ifPDFTeX
  \usepackage[T1]{fontenc}
  \usepackage[utf8]{inputenc}
  \usepackage{textcomp} % provide euro and other symbols
\else % if luatex or xetex
  \usepackage{unicode-math} % this also loads fontspec
  \defaultfontfeatures{Scale=MatchLowercase}
  \defaultfontfeatures[\rmfamily]{Ligatures=TeX,Scale=1}
\fi
\usepackage{lmodern}
\ifPDFTeX\else
  % xetex/luatex font selection
\fi
% Use upquote if available, for straight quotes in verbatim environments
\IfFileExists{upquote.sty}{\usepackage{upquote}}{}
\IfFileExists{microtype.sty}{% use microtype if available
  \usepackage[]{microtype}
  \UseMicrotypeSet[protrusion]{basicmath} % disable protrusion for tt fonts
}{}
\makeatletter
\@ifundefined{KOMAClassName}{% if non-KOMA class
  \IfFileExists{parskip.sty}{%
    \usepackage{parskip}
  }{% else
    \setlength{\parindent}{0pt}
    \setlength{\parskip}{6pt plus 2pt minus 1pt}}
}{% if KOMA class
  \KOMAoptions{parskip=half}}
\makeatother
% Make \paragraph and \subparagraph free-standing
\makeatletter
\ifx\paragraph\undefined\else
  \let\oldparagraph\paragraph
  \renewcommand{\paragraph}{
    \@ifstar
      \xxxParagraphStar
      \xxxParagraphNoStar
  }
  \newcommand{\xxxParagraphStar}[1]{\oldparagraph*{#1}\mbox{}}
  \newcommand{\xxxParagraphNoStar}[1]{\oldparagraph{#1}\mbox{}}
\fi
\ifx\subparagraph\undefined\else
  \let\oldsubparagraph\subparagraph
  \renewcommand{\subparagraph}{
    \@ifstar
      \xxxSubParagraphStar
      \xxxSubParagraphNoStar
  }
  \newcommand{\xxxSubParagraphStar}[1]{\oldsubparagraph*{#1}\mbox{}}
  \newcommand{\xxxSubParagraphNoStar}[1]{\oldsubparagraph{#1}\mbox{}}
\fi
\makeatother


\usepackage{longtable,booktabs,array}
\usepackage{calc} % for calculating minipage widths
% Correct order of tables after \paragraph or \subparagraph
\usepackage{etoolbox}
\makeatletter
\patchcmd\longtable{\par}{\if@noskipsec\mbox{}\fi\par}{}{}
\makeatother
% Allow footnotes in longtable head/foot
\IfFileExists{footnotehyper.sty}{\usepackage{footnotehyper}}{\usepackage{footnote}}
\makesavenoteenv{longtable}
\usepackage{graphicx}
\makeatletter
\newsavebox\pandoc@box
\newcommand*\pandocbounded[1]{% scales image to fit in text height/width
  \sbox\pandoc@box{#1}%
  \Gscale@div\@tempa{\textheight}{\dimexpr\ht\pandoc@box+\dp\pandoc@box\relax}%
  \Gscale@div\@tempb{\linewidth}{\wd\pandoc@box}%
  \ifdim\@tempb\p@<\@tempa\p@\let\@tempa\@tempb\fi% select the smaller of both
  \ifdim\@tempa\p@<\p@\scalebox{\@tempa}{\usebox\pandoc@box}%
  \else\usebox{\pandoc@box}%
  \fi%
}
% Set default figure placement to htbp
\def\fps@figure{htbp}
\makeatother


% definitions for citeproc citations
\NewDocumentCommand\citeproctext{}{}
\NewDocumentCommand\citeproc{mm}{%
  \begingroup\def\citeproctext{#2}\cite{#1}\endgroup}
\makeatletter
 % allow citations to break across lines
 \let\@cite@ofmt\@firstofone
 % avoid brackets around text for \cite:
 \def\@biblabel#1{}
 \def\@cite#1#2{{#1\if@tempswa , #2\fi}}
\makeatother
\newlength{\cslhangindent}
\setlength{\cslhangindent}{1.5em}
\newlength{\csllabelwidth}
\setlength{\csllabelwidth}{3em}
\newenvironment{CSLReferences}[2] % #1 hanging-indent, #2 entry-spacing
 {\begin{list}{}{%
  \setlength{\itemindent}{0pt}
  \setlength{\leftmargin}{0pt}
  \setlength{\parsep}{0pt}
  % turn on hanging indent if param 1 is 1
  \ifodd #1
   \setlength{\leftmargin}{\cslhangindent}
   \setlength{\itemindent}{-1\cslhangindent}
  \fi
  % set entry spacing
  \setlength{\itemsep}{#2\baselineskip}}}
 {\end{list}}
\usepackage{calc}
\newcommand{\CSLBlock}[1]{\hfill\break\parbox[t]{\linewidth}{\strut\ignorespaces#1\strut}}
\newcommand{\CSLLeftMargin}[1]{\parbox[t]{\csllabelwidth}{\strut#1\strut}}
\newcommand{\CSLRightInline}[1]{\parbox[t]{\linewidth - \csllabelwidth}{\strut#1\strut}}
\newcommand{\CSLIndent}[1]{\hspace{\cslhangindent}#1}



\setlength{\emergencystretch}{3em} % prevent overfull lines

\providecommand{\tightlist}{%
  \setlength{\itemsep}{0pt}\setlength{\parskip}{0pt}}



 


\KOMAoption{captions}{tableheading}
\makeatletter
\@ifpackageloaded{caption}{}{\usepackage{caption}}
\AtBeginDocument{%
\ifdefined\contentsname
  \renewcommand*\contentsname{Table of contents}
\else
  \newcommand\contentsname{Table of contents}
\fi
\ifdefined\listfigurename
  \renewcommand*\listfigurename{List of Figures}
\else
  \newcommand\listfigurename{List of Figures}
\fi
\ifdefined\listtablename
  \renewcommand*\listtablename{List of Tables}
\else
  \newcommand\listtablename{List of Tables}
\fi
\ifdefined\figurename
  \renewcommand*\figurename{Figure}
\else
  \newcommand\figurename{Figure}
\fi
\ifdefined\tablename
  \renewcommand*\tablename{Table}
\else
  \newcommand\tablename{Table}
\fi
}
\@ifpackageloaded{float}{}{\usepackage{float}}
\floatstyle{ruled}
\@ifundefined{c@chapter}{\newfloat{codelisting}{h}{lop}}{\newfloat{codelisting}{h}{lop}[chapter]}
\floatname{codelisting}{Listing}
\newcommand*\listoflistings{\listof{codelisting}{List of Listings}}
\makeatother
\makeatletter
\makeatother
\makeatletter
\@ifpackageloaded{caption}{}{\usepackage{caption}}
\@ifpackageloaded{subcaption}{}{\usepackage{subcaption}}
\makeatother
\usepackage{bookmark}
\IfFileExists{xurl.sty}{\usepackage{xurl}}{} % add URL line breaks if available
\urlstyle{same}
\hypersetup{
  pdftitle={Stroke reorganises the transient dynamical structure of human gait},
  colorlinks=true,
  linkcolor={blue},
  filecolor={Maroon},
  citecolor={Blue},
  urlcolor={Blue},
  pdfcreator={LaTeX via pandoc}}


\title{Stroke reorganises the transient dynamical structure of human
gait}
\author{Michael Hess \and Gordon Berman}
\date{}
\begin{document}
\maketitle


\subsection{Abstract}\label{abstract}

Healthy walking requires flexible, multi-joint coordination that can
absorb perturbations across a range of frequencies and directions while
maintaining rhythmic locomotion. Classical gait stability measures --
Floquet multipliers and Lyapunov exponents -- reduce this rich
coordination to a single number summarising how fast perturbations grow
or decay, discarding information about \emph{which} perturbation
directions matter, \emph{how} the system responds at different
frequencies, and \emph{how} the two limbs couple to each other. Here we
fit a simple linear model (delay-embedded dynamic mode decomposition) to
bilateral hip, knee, and ankle kinematics from 299 treadmill walking
trials and show that the structure of this model -- specifically, how
far it departs from a ``well-behaved'' (normal) operator -- provides a
multi-dimensional portrait of locomotor coordination that distinguishes
able-bodied, high-functioning stroke, and low-functioning stroke gait in
ways that conventional metrics cannot. Able-bodied operators are highly
ill-conditioned, enabling the model to respond broadly across stride
harmonics and maintain symmetric inter-limb coupling regardless of
walking speed. High-functioning stroke operators concentrate their
response sharply at the stride frequency, achieving strong resonance but
sacrificing frequency breadth. Low-functioning stroke operators are
nearly well-conditioned, producing weak responses across all
frequencies. Both stroke groups show speed-dependent non-normality
(absent in able-bodied gait) and increased inter-limb coupling asymmetry
concentrated around the non-paretic knee. These signatures reveal that
stroke does not merely slow gait or increase variability: it
restructures the dynamical coordination strategy in group-specific ways,
suggesting that rehabilitation approaches targeting dynamical richness
may complement conventional kinematic goals.

\subsection{Introduction}\label{introduction}

Walking looks simple but requires remarkably sophisticated coordination.
Each stride involves dozens of muscles acting across multiple joints on
both legs, timed to maintain balance, absorb ground reaction forces, and
propel the body forward. When this coordination is disrupted by
neurological injury such as stroke, the consequences are visible not
just as slower walking or asymmetric steps, but as a qualitative change
in how the whole system responds to the moment-to-moment perturbations
-- from uneven ground, shifting attention, or the body's own variability
-- that healthy walkers handle without conscious effort.

How should we measure this coordination? The dominant approach treats
walking as a periodic orbit and asks a single question: do perturbations
shrink or grow? Floquet multipliers answer this per stride cycle, but
require strict periodicity and very large datasets -- Revzen and
Guckenheimer (2008) needed over 10,000 strides from cockroach running
data just to obtain stable estimates, orders of magnitude more than the
\textasciitilde15--50 strides in a typical human trial. They also demand
a reliable phase estimate, itself a hard problem (Revzen and
Guckenheimer 2012; Bruijn et al. 2013). The maximum Lyapunov exponent
(\(\lambda_S\)) relaxes the periodicity assumption and has emerged as
the best-validated predictor of fall risk (Dingwell and Cusumano 2000;
Bruijn et al. 2013), but it provides only a scalar: how fast do nearby
trajectories pull apart? It says nothing about \emph{which directions}
are unstable, \emph{at what frequencies} the system can respond, or
\emph{how the two limbs talk to each other}. These are precisely the
questions that matter for understanding locomotor coordination after
stroke.

Both measures also share a deeper conceptual issue: they assume the
system is autonomous -- evolving under its own dynamics without external
input. Walking is not autonomous. It is continuously driven by
descending neural commands, modulated by sensory feedback, and
stabilised by reflexes (Ting and McKay 2007). The common finding that
\(\lambda_S\) is positive (local divergence) despite stable walking is a
direct consequence: the system is being simultaneously driven and
controlled, and Lyapunov analysis of the autonomous reconstruction picks
up the driving, not a genuine instability (Dingwell and Kang 2007).

We take a different approach. Rather than asking a scalar stability
question, we fit a simple linear model to the walking data and then
examine the \emph{structure of that model} to learn about the underlying
coordination. Specifically, we use dynamic mode decomposition (DMD) -- a
data-driven technique that finds the best-fit linear operator \(A\)
predicting the next time step from the current one (Schmid 2010; Tu et
al. 2014) -- applied to delay-embedded bilateral joint kinematics. Delay
embedding stacks successive time points into a higher-dimensional state
vector, allowing the linear model to capture temporal dependencies that
a single snapshot cannot (Takens 1981; Brunton and Kutz 2019).

Like Floquet and Lyapunov analysis, DMD fits an autonomous model. But
rather than reducing the result to eigenvalues (which determine long-run
stability), we examine the model's \emph{non-normality}: how far the
fitted operator departs from having orthogonal eigenvectors. This
departure has concrete physical meaning. A ``well-behaved'' (normal)
operator can only respond to perturbations through its individual modes,
each decaying independently. A non-normal operator can transiently
amplify specific perturbation directions, respond at frequencies where
no single mode is resonant, and couple inputs and outputs in
directionally specific ways (Trefethen and Embree 2005; Henrici 1962).
These are not abstract mathematical properties -- they map directly onto
the locomotor questions we care about: How flexibly can the system
respond to perturbations across frequencies? How do the two limbs
influence each other? And does the system's coordination structure
depend on walking speed, or is it robust across operating conditions?

Non-normality has previously been studied in the DMD literature
primarily as a \emph{numerical problem}: Wu, Brunton, and Revzen (2021)
showed that when DMD is applied to oscillatory data, the underlying
operator is inherently non-normal, leading to ill-conditioning and
noise-amplified errors. We invert this perspective. Rather than treating
non-normality as a nuisance, we show that its magnitude and structure --
recovered from walking kinematics -- are biologically informative
signatures of locomotor coordination. Where the fitted linear model must
strain hardest to approximate the true, nonlinear, driven system,
non-normality is highest. The pattern of that strain tells us about the
coordination it is trying to capture.

Post-stroke gait provides a compelling test case. The clinical picture
includes slower walking speeds, inter-limb asymmetry, and elevated
stride-to-stride variability (Bruijn et al. 2013), but these
descriptions remain at the level of symptoms. A deeper question is
whether stroke fundamentally reorganises \emph{how} the locomotor system
coordinates itself -- not just how fast or how symmetrically it walks,
but the structure of its perturbation responses, its frequency content,
and its inter-limb coupling. We test this hypothesis using 299 treadmill
walking trials from the open dataset of Winner et al. (2023), which
includes able-bodied controls and high- and low-functioning stroke
survivors. Our approach complements their LSTM-based ``gait signatures''
framework: where that work used black-box neural networks to identify
individual-specific dynamics, we use an explicitly analysable linear
operator to extract interpretable coordination signatures.

\subsection{Methods}\label{methods}

\subsubsection{Dataset}\label{dataset}

We analyse 299 treadmill walking trials drawn from the open dataset of
Winner et al. (2023). Each trial comprises 1500 time points (15 s at 100
Hz) of six kinematic channels: bilateral hip, knee, and ankle
sagittal-plane joint angles. Trials are partitioned by metadata into
able-bodied (AB, \(n = 180\); speed range 30--220 cm/s),
high-functioning stroke (HF, \(n = 65\); 40--127 cm/s), and
low-functioning stroke (LF, \(n = 54\); 22--60 cm/s). All channels are
z-scored within each trial prior to analysis, so that operator
properties reflect dynamical organisation rather than absolute amplitude
differences.

\subsubsection{Delay-Embedded Hankel
DMD}\label{delay-embedded-hankel-dmd}

We construct a delay-embedded state by stacking \(\tau = 10\) lags of
the 6-channel observation, yielding a \(66\)-dimensional augmented state
vector at each time step (Takens 1981). This embedding lag spans
approximately one stride period at typical walking cadences (100 Hz
sampling), ensuring that the augmented space captures at least one full
cycle of the dominant oscillation. The resulting Hankel matrix
\(H \in \mathbb{R}^{66 \times (T - \tau)}\) is split into input
(\(X = H_{:,1:end-1}\)) and output (\(Y = H_{:,2:end}\)) matrices. We
compute the truncated SVD \(X = U \Sigma V^*\), retaining the \(r\)
singular values exceeding \(10^{-6}\) times the leading singular value
(typically retaining \textgreater99\% of the variance), and solve for
the reduced operator \(\tilde{A} = U_r^* Y V_r \Sigma_r^{-1}\) following
the standard projected DMD formulation (Tu et al. 2014; Schmid 2010). To
enforce stability, eigenvalues of \(\tilde{A}\) with modulus exceeding
unity are projected onto the unit circle. All subsequent analyses
operate on this reduced operator.

\subsubsection{Dynamical Quantities}\label{dynamical-quantities}

\textbf{Henrici departure from normality.} We quantify global
non-normality as \(\eta = \|A^* A - A A^*\|_2 \,/\, \|A\|_2^2\), where
\(\|\cdot\|_2\) denotes the spectral (operator 2-) norm. A normal
operator satisfies \(\eta = 0\); larger values indicate stronger
non-orthogonality of eigenvectors and correspondingly greater potential
for transient amplification (Henrici 1962; Trefethen and Embree 2005).

\textbf{Harmonic resolvent gain.} The stride frequency \(f_s\) is
identified as the dominant peak (above 0.3 Hz) in the resolvent gain
\(\sigma_{\max}^{-1}(e^{i\theta}I - A)\) swept around the unit circle.
Harmonic resolvent gains are then evaluated at
\(\|(e^{i 2\pi h f_s / f_{\text{samp}}} I - A)^{-1}\|\) for harmonics
\(h = 1, \ldots, 5\), quantifying how strongly the operator amplifies
periodic forcing at each multiple of the stride frequency.

\textbf{Spatial resolvent asymmetry.} The full resolvent at the stride
frequency is mapped back to the original \(66\)-dimensional
(delay-embedded) joint space via
\(R = U \,(z I - \tilde{A})^{-1}\, U^*\). We extract sub-blocks
corresponding to paretic\(\to\)non-paretic and non-paretic\(\to\)paretic
coupling (across all delay lags) and compute their operator norms. The
asymmetry index is \(\max(G_{12}, G_{21}) / \min(G_{12}, G_{21})\),
where \(G_{12}\) and \(G_{21}\) are the respective sub-block norms.
Signed log-ratios of pairwise joint gains quantify the directionality of
inter-joint coupling.

\textbf{Peak Q-factor.} The quality factor at the stride frequency is
\(Q = f_s / \Delta f\), where \(\Delta f\) is the half-power bandwidth
(frequency interval over which the resolvent gain exceeds
\(\text{peak} / \sqrt{2}\)). Higher Q indicates narrower frequency
tuning.

\textbf{Operator condition number.} The condition number
\(\kappa(A) = \sigma_{\max}(A) / \sigma_{\min}(A)\) quantifies the
ill-conditioning of the operator. High \(\kappa\) implies highly
non-orthogonal eigenvectors and determines the extent of
\emph{pseudospectral spreading}: how far the resolvent gain extends away
from eigenvalue locations on the unit circle. For a normal operator,
\(\kappa = 1\) and the resolvent is determined entirely by eigenvalue
distances; for a non-normal operator, the pseudospectral
\(\varepsilon\)-contours broaden with \(\kappa\), elevating gain at
frequencies where no eigenvalue sits nearby (Trefethen and Embree 2005).

\textbf{Eigenvalue--stride distance.} We also report
\(\min_j |\lambda_j - e^{i 2\pi f_s / f_{\text{samp}}}|\), the distance
from the nearest eigenvalue of \(\tilde{A}\) to the stride-frequency
point on the unit circle. This quantity controls the peak resolvent gain
at the fundamental: tighter proximity yields larger absolute gain,
independent of non-normality.

All quantities are computed trial-by-trial. Speed-matched comparisons
cap AB trials at the maximum LF walking speed (60 cm/s) to control for
the confound that AB subjects walk faster on average.

\subsubsection{Statistical Inference}\label{statistical-inference}

Because each participant contributes multiple trials at different
walking speeds, inferential comparisons use a hierarchical bootstrap
rather than treating trials as independent replicates. Within each
group, subjects are resampled with replacement, and then that subject's
trials are resampled with replacement within the bootstrap draw. This
preserves the nested subject/trial structure while accommodating unequal
trial counts per subject. For between-group contrasts we bootstrap
differences in medians; for speed-dependence we bootstrap Pearson
correlations; and for directional quantities expressed as log-ratios we
bootstrap the mean log-ratio. We report percentile 95\% confidence
intervals and two-sided bootstrap p-values computed from the fraction of
bootstrap samples on either side of zero. Unless otherwise noted,
cross-group inference is reported on the speed-matched subset.

\subsection{Results}\label{results}

\textsubscript{Source:
\href{https://Michaelhess17.github.io/Synology/index.qmd.html}{Article
Notebook}}

\subsubsection{Figure 1: Non-Normality Tracks Nonlinearity and Timescale
Separation in Simulated
Systems}\label{figure-1-non-normality-tracks-nonlinearity-and-timescale-separation-in-simulated-systems}

\textsubscript{Source:
\href{https://Michaelhess17.github.io/Synology/index.qmd.html}{Article
Notebook}}

\begin{verbatim}

Saved: figures/fig1_simulated_benchmarks.svg
\end{verbatim}

\begin{figure}[H]

\centering{

\pandocbounded{\includegraphics[keepaspectratio]{index_files/mediabag/index_files/figure-pdf/fig-simulated-output-2.pdf}}

}

\caption{\label{fig-simulated}Henrici departure in simulated
oscillators. (A) Stuart--Landau: Henrici rises with μ as the limit cycle
amplitude increases. (B) FitzHugh-Nagumo: Henrici rises as ε decreases
and timescale separation strengthens. (C) Stuart--Landau with isochron
shear: Henrici rises with shear parameter c.~(D) Joint heatmap over (μ,
c).}

\end{figure}%

\textsubscript{Source:
\href{https://Michaelhess17.github.io/Synology/index.qmd.html}{Article
Notebook}}

Simulated benchmark oscillators establish that Henrici departure in a
delay-embedded DMD operator rises with three properties of the
underlying nonlinear flow: (A) limit-cycle amplitude (Stuart--Landau,
increasing \(\mu\)), (B) timescale separation (FitzHugh--Nagumo,
decreasing \(\varepsilon\)), and (C) isochron shear (Stuart--Landau with
amplitude-dependent frequency, increasing \(c\)). Panel D confirms that
these effects combine monotonically. In each case the linear surrogate
encodes the nonlinear geometry through nearly parallel eigenvectors
whose large positive and negative contributions cancel on the cycle but
amplify deviations off it. These simulations supply the mechanistic
interpretation used throughout: elevated Henrici departure indexes
richer transient structure, reflecting the complexity of the underlying
nonlinear dynamics that the linear operator must approximate.

\subsubsection{Figure 2: Speed-Dependent Non-Normality Distinguishes
Stroke from Able-Bodied
Gait}\label{figure-2-speed-dependent-non-normality-distinguishes-stroke-from-able-bodied-gait}

\textsubscript{Source:
\href{https://Michaelhess17.github.io/Synology/index.qmd.html}{Article
Notebook}}

\begin{figure}[H]

\centering{

\pandocbounded{\includegraphics[keepaspectratio]{index_files/mediabag/index_files/figure-pdf/fig-nonnormality-output-1.pdf}}

}

\caption{\label{fig-nonnormality}Henrici departure vs walking speed for
299 gait trials. AB = able-bodied (steelblue), HF = high-functioning
stroke (orange), LF = low-functioning stroke (crimson). Lines connect
trials from the same subject across speeds. The strip chart (right)
shows speed-matched comparisons.}

\end{figure}%

\textsubscript{Source:
\href{https://Michaelhess17.github.io/Synology/index.qmd.html}{Article
Notebook}}

\begin{figure}[H]

\centering{

\pandocbounded{\includegraphics[keepaspectratio]{index_files/mediabag/index_files/figure-pdf/fig-condnum-output-1.pdf}}

}

\caption{\label{fig-condnum}Operator condition number κ(A) by group
(speed-matched, log scale). Higher conditioning indicates more
ill-conditioned eigenvector geometry---a prerequisite for pseudospectral
spreading of resolvent gain away from eigenvalue locations.}

\end{figure}%

\textsubscript{Source:
\href{https://Michaelhess17.github.io/Synology/index.qmd.html}{Article
Notebook}}

Able-bodied gait maintains high non-normality (median \(\eta\) = 0.96,
IQR {[}0.95, 0.99{]}) that is not significantly correlated with walking
speed (\(r\) = 0.11, 95\% CI {[}-0.04, 0.26{]}, \(p\) = 0.14;
hierarchical bootstrap). Both stroke groups show significantly stronger
speed-dependence: high-functioning stroke achieves the highest peak
non-normality (median \(\eta\) = 1.0, \(r\) = 0.35, CI {[}0.1, 0.56{]},
\(p\) = 0.012), while low-functioning stroke exhibits the lowest overall
values (median \(\eta\) = 0.89, \(r\) = 0.47, CI {[}0.12, 0.75{]}, \(p\)
= 0.016). Speed-matched comparisons confirm that HF non-normality
exceeds both AB (\(\Delta\)median = 0.047, \(p\) = 0.018) and LF
(\(\Delta\)median = 0.106, \(p\) = 0.005), while the LF--AB contrast
trends negative but remains uncertain (\(\Delta\)median = -0.059, \(p\)
= 0.063).

A second metric, the operator condition number \(\kappa(A)\), measures
something related but distinct: how broadly the model can spread its
response across frequencies beyond the locations of its individual
modes. A high condition number means the model's response ``leaks''
widely -- perturbations at one frequency can be amplified at
neighbouring frequencies even where no mode sits. A low condition number
confines the response tightly around each mode. The groups are sharply
separated: AB operators have extremely high conditioning (median
\(\kappa \approx\) 275), HF moderate (\(\kappa \approx\) 1475), and LF
nearly minimal (\(\kappa \approx\) 78). These differences are
significant on the log scale (AB vs LF \(\Delta\)median log \(\kappa\) =
-1.3, \(p\) = 0.006; AB vs HF \(\Delta\)median log \(\kappa\) = 1.7,
\(p\) = 0.552). AB operators also retain higher SVD rank (median 45 vs
37 HF and 36 LF out of 66 possible dimensions), indicating that the
able-bodied coordination occupies a higher-dimensional subspace. As we
show in the next section, this three-way gradient in conditioning --
high (AB), moderate (HF), low (LF) -- explains why the groups' frequency
response profiles take such different shapes despite comparable or even
elevated non-normality in HF.

\subsubsection{Figure 3: Harmonic Resolvent Gain Reveals Stroke-Related
Collapse of High-Frequency
Coupling}\label{figure-3-harmonic-resolvent-gain-reveals-stroke-related-collapse-of-high-frequency-coupling}

\textsubscript{Source:
\href{https://Michaelhess17.github.io/Synology/index.qmd.html}{Article
Notebook}}

\begin{figure}[H]

\centering{

\pandocbounded{\includegraphics[keepaspectratio]{index_files/mediabag/index_files/figure-pdf/fig-harmonic-all-output-1.pdf}}

}

\caption{\label{fig-harmonic-all}Absolute and normalised harmonic
resolvent gain. (A) Absolute gain at harmonics 1--5. (B) Gain normalised
by the fundamental, revealing fall-off shape independent of overall gain
level.}

\end{figure}%

\textsubscript{Source:
\href{https://Michaelhess17.github.io/Synology/index.qmd.html}{Article
Notebook}}

\begin{figure}[H]

\centering{

\pandocbounded{\includegraphics[keepaspectratio]{index_files/mediabag/index_files/figure-pdf/fig-harmonic-speed-output-1.pdf}}

}

\caption{\label{fig-harmonic-speed}Normalised harmonic gain faceted by
walking speed bin (blue = slow, red = fast). Speed modulates the gain
profile within each group differently, with AB showing the largest
speed-dependent change.}

\end{figure}%

\textsubscript{Source:
\href{https://Michaelhess17.github.io/Synology/index.qmd.html}{Article
Notebook}}

\begin{figure}[H]

\centering{

\pandocbounded{\includegraphics[keepaspectratio]{index_files/mediabag/index_files/figure-pdf/fig-harmonic-matched-output-1.pdf}}

}

\caption{\label{fig-harmonic-matched}Speed-matched harmonic gain
comparison (AB capped at the maximum LF speed). Group differences in the
fall-off shape persist after controlling for speed.}

\end{figure}%

\textsubscript{Source:
\href{https://Michaelhess17.github.io/Synology/index.qmd.html}{Article
Notebook}}

These profiles show how broadly each group's model can respond across
multiples of the stride frequency -- a proxy for how rich and
multi-timescale the underlying coordination is. The pattern is more
nuanced than ``stroke suppresses higher harmonics.''

To quantify the fall-off, we fit a spectral decay rate (slope of
log-normalised gain vs.~harmonic number) to each trial. AB gait decays
at a moderate rate (median slope -1.19), LF is similar (-1.11), and HF
decays significantly faster (-1.86). Hierarchical bootstrap confirms HF
falls off faster than AB (\(\Delta\)median = -0.9, \(p\) = 0.001); LF is
not significantly different from AB (\(\Delta\)median = -0.16, \(p\) =
0.232), and the HF--LF contrast trends in the same direction but is
noisier (\(p\) = 0.208).

This result initially seems paradoxical: HF has the highest
non-normality (Figure 2) yet the steepest fall-off -- shouldn't a more
``complex'' operator respond more broadly? The resolution involves two
distinct mechanisms. First, \emph{how closely the model's internal
oscillation matches the stride frequency} determines the absolute gain
at the fundamental. HF operators place their nearest mode very close to
the stride frequency (median distance 0.00069 vs.~0.0012 for AB),
producing the highest absolute gain at h1 (median 1.9 \(\times 10^7\)
vs.~5.3 \(\times 10^6\) AB; bootstrap \(\Delta\)median = 1.5
\(\times 10^7\), \(p\) \textless{} 0.001). Second, \emph{how
ill-conditioned the operator is} (the condition number \(\kappa\),
reported in Figure 2) determines how broadly the model's response
spreads away from its modes to neighbouring frequencies. AB's extreme
ill-conditioning (\(\kappa \approx\) 18,700) lets the response ``leak''
to harmonics that no single mode directly supports. HF's moderate
conditioning (\(\kappa \approx\) 1,200) limits this spreading.

Crucially, the steep normalised fall-off in HF does \emph{not} mean
higher harmonics are suppressed in absolute terms. At h2--h5 the
absolute gains are comparable across groups or even highest in LF (h2:
AB 0.7M, HF 0.9M, LF 1.6M). Rather, HF's fundamental is so strongly
amplified that everything else looks small by comparison.

The biological picture: able-bodied gait uses a high-dimensional,
broadly responsive model -- the system can handle perturbations at many
frequencies. HF gait is sharply tuned to the stride frequency, achieving
strong resonance but sacrificing breadth. LF gait responds weakly across
all frequencies, with no dominant resonance.

\subsubsection{Figure 4: Spatial Resolvent Asymmetry, Q-Factor, and
Knee-Centred Coupling Summaries Differentiate
Groups}\label{figure-4-spatial-resolvent-asymmetry-q-factor-and-knee-centred-coupling-summaries-differentiate-groups}

\textsubscript{Source:
\href{https://Michaelhess17.github.io/Synology/index.qmd.html}{Article
Notebook}}

\begin{figure}[H]

\centering{

\pandocbounded{\includegraphics[keepaspectratio]{index_files/mediabag/index_files/figure-pdf/fig-novel-strips-output-1.pdf}}

}

\caption{\label{fig-novel-strips}Speed-matched strip charts for the
coarse operator metrics. (A) Spatial resolvent asymmetry: ratio of
cross-body resolvent norms, max/min. (B) Peak Q-factor at stride
frequency. (C) Mean paretic→non-paretic resolvent gain. (D) Mean
non-paretic→paretic resolvent gain. Lines show median; error bars show
IQR.}

\end{figure}%

\textsubscript{Source:
\href{https://Michaelhess17.github.io/Synology/index.qmd.html}{Article
Notebook}}

\begin{figure}[H]

\centering{

\pandocbounded{\includegraphics[keepaspectratio]{index_files/mediabag/index_files/figure-pdf/fig-joint-summary-output-1.pdf}}

}

\caption{\label{fig-joint-summary}Speed-matched summary metrics
replacing the full joint-level heatmaps. (E) Knee-target asymmetry,
defined as log(hip→non-paretic knee / hip→paretic knee). Positive values
indicate that hip-to-knee coupling is concentrated on the non-paretic
knee. (F) Non-paretic-knee directional bias, defined as the mean
log{[}(hip→non-paretic knee)/(non-paretic knee→hip){]} across both hips.
Positive values indicate preferential drive into the non-paretic knee.
(G) Global pairwise asymmetry, defined as the mean absolute signed
log-ratio across all joint pairs.}

\end{figure}%

\textsubscript{Source:
\href{https://Michaelhess17.github.io/Synology/index.qmd.html}{Article
Notebook}}

In healthy walking, the two legs influence each other roughly
symmetrically -- perturbations on one side propagate to the other with
similar strength in both directions. Stroke disrupts this balance. We
quantify inter-limb coupling asymmetry as the ratio between the stronger
and weaker cross-body influence directions (a ratio of 1 means perfect
symmetry). This ratio rises monotonically from AB (median 1.23) through
HF (1.69) to LF (1.89), indicating progressively one-sided coupling
between the legs. Hierarchical bootstrap confirms larger asymmetry in
both stroke groups relative to AB (HF--AB \(\Delta\)median = 0.47, 95\%
CI {[}0.24, 0.92{]}, \(p\) = 0.002; LF--AB \(\Delta\) = 0.66, {[}0.28,
1.75{]}, \(p\) \textless{} 0.001). Moreover, in LF gait the dominant
direction of influence runs from the paretic to the non-paretic side
(mean log-ratio = 0.6, 95\% CI {[}0.28, 0.9{]}, \(p\) \textless{}
0.001), with HF trending similarly but less certain (0.21, {[}-0.01,
0.45{]}, \(p\) = 0.062). In other words, the impaired limb's dynamics
propagate forcefully across the body, and the intact side must absorb
and accommodate those perturbations.

How narrowly tuned is each group's coordination to the stride frequency?
The Q-factor measures this: higher Q means the system responds strongly
at stride frequency but poorly at neighbouring frequencies. HF gait is
the most sharply tuned (median Q = 69.6 vs.~AB 14.3 and LF 31.3),
consistent with a coordination strategy that locks onto the stride
rhythm but sacrifices adaptability to perturbations at other
frequencies. Hierarchical bootstrap confirms higher Q in HF than AB
(\(\Delta\)median = 55.3, 95\% CI {[}24.6, 93.5{]}, \(p\) \textless{}
0.001) and in LF than AB (\(\Delta\) = 17.1, {[}2.5, 59.6{]}, \(p\) =
0.012), while the HF--LF contrast is less certain (\(p\) = 0.176). LF
gait thus combines high spatial asymmetry with moderately narrow tuning
-- a coordination pattern that is both directionally biased and somewhat
rigid.

To localise where in the kinematic chain inter-limb asymmetry is
concentrated, we examine hip-to-knee coupling. In AB gait, the hips
drive both knees roughly equally. In both stroke groups, this coupling
shifts toward the non-paretic knee: the log-ratio of
hip-to-non-paretic-knee vs.~hip-to-paretic-knee coupling is near zero in
AB but positive in stroke (HF--AB \(\Delta\)median = 0.23, 95\% CI
{[}0.01, 0.65{]}, \(p\) = 0.034; LF--AB \(\Delta\) = 0.79, {[}0.12,
1.39{]}, \(p\) = 0.016). Both stroke groups also show strong directional
bias \emph{into} the non-paretic knee from the hips (HF: mean = 1.09,
95\% CI {[}0.82, 1.36{]}, \(p\) \textless{} 0.001; LF: mean = 1.19,
{[}0.87, 1.48{]}, \(p\) \textless{} 0.001), with LF exceeding AB
(\(\Delta\)median = 0.83, {[}0.18, 1.33{]}, \(p\) = 0.012). A broader
summary of all pairwise joint asymmetries trends upward in LF relative
to AB (\(\Delta\)median = 0.33, 95\% CI {[}-0.02, 0.57{]}, \(p\) =
0.056), but the clearest finding is that stroke-related coupling
asymmetry is concentrated specifically around the non-paretic knee
rather than being diffusely spread across all joints. The intact leg
does not simply ``compensate'' in a general sense -- the model reveals a
specific pathway through which the non-paretic knee absorbs and
redistributes perturbations originating from the impaired side.

\subsection{Discussion}\label{discussion}

Three operator-level signatures emerge from delay-embedded DMD of 299
gait trials, collectively revealing that stroke reorganises--rather than
merely degrades--the transient dynamical structure of locomotion.

\textbf{Speed-dependent non-normality.} Able-bodied gait maintains a
stable, speed-invariant level of non-normality, indicating that the
transient complexity of the locomotor operator is robust across
operating points. Both stroke groups show elevated speed-dependence
(\(r = 0.35\) HF, \(r = 0.47\) LF vs.~\(r = 0.11\) AB), meaning that
their capacity for transient amplification is fragile--contingent on
walking at particular speeds. As discussed in the Introduction, all
three analysis frameworks--Floquet multipliers, Lyapunov exponents, and
DMD--fit autonomous models to a system that is continuously driven and
controlled (Dingwell and Kang 2007; Ting and McKay 2007). The maximum
Lyapunov exponent \(\lambda_S\) remains the best-validated predictor of
fall risk (Bruijn et al. 2013), but as a scalar divergence rate of an
autonomous reconstruction it cannot distinguish intrinsic instability
from neural driving, nor can it capture the directional or
frequency-dependent structure of transient dynamics. The Henrici
departure, by indexing how severely the autonomous linear model must
distort its eigenvector geometry to fit the driven, nonlinear data,
captures a complementary dimension of locomotor organisation. Where
\(\lambda_S\) summarises the dominant instability direction, the
resolvent structure maps the full landscape of transient amplification
across frequencies and spatial channels. The speed-dependence of this
structure in stroke gait--absent in able-bodied gait--suggests that the
model-reality gap (i.e., how ``non-autonomous'' or ``nonlinear'' the
true system is relative to the linear surrogate) varies with operating
point in impaired but not in healthy locomotion.

\textbf{Spectral organisation: eigenvalue proximity vs.~pseudospectral
spreading.} The harmonic resolvent profiles initially appear
paradoxical: HF has the highest Henrici departure yet the steepest
normalised harmonic fall-off, while LF has the lowest non-normality yet
a normalised profile similar to AB. The resolution lies in recognising
that two distinct mechanisms shape the resolvent--\emph{eigenvalue
proximity} to the stride frequency (which sets absolute gain at each
harmonic) and \emph{pseudospectral spreading} via the operator's
condition number (which determines how broadly gain extends away from
eigenvalue locations). HF operators park an eigenvalue very close to the
stride fundamental, producing massive h1 gain. But HF's condition number
(\(\kappa \approx 1{,}200\)) is an order of magnitude below AB's
(\(\kappa \approx 18{,}700\)), so the pseudospectral halo that would
spread this gain to neighbouring harmonics is narrower. The steep
normalised fall-off in HF therefore reflects disproportionate inflation
of the fundamental, not suppression of higher harmonics--absolute gains
at h2--h5 are comparable to or exceed AB's. Conversely, LF operators are
nearly well-conditioned (\(\kappa \approx 46\)), meaning almost no
pseudospectral spreading; their flat normalised profile arises simply
because no single eigenvalue dominates the stride frequency. AB combines
extreme ill-conditioning with moderate eigenvalue--stride coupling to
produce both high absolute gain and broad bandwidth--the pseudospectral
effects sustain resolvent gain at frequencies far from any eigenvalue.
From the perspective of Wu, Brunton, and Revzen (2021), the condition
number gradient across groups reflects how differently the Koopman
operator's non-normal structure is distorted by neurological impairment.
The non-normality they identify as a challenge for DMD convergence here
becomes a graded diagnostic: AB dynamics strain the autonomous linear
model broadly across the spectrum; HF dynamics strain it acutely at one
frequency; LF dynamics strain it least of all.

\textbf{Asymmetric coupling focused on the non-paretic knee.} The
directionality of inter-limb coupling provides perhaps the most
clinically informative signature. At the coarse inter-limb level, LF
shows clear paretic-dominant cross-body drive, with HF trending in the
same direction. At finer joint resolution, however, the most robust
effect is even more specific: hip-knee coupling becomes concentrated on
the non-paretic knee, and both stroke groups show strong directional
bias into that knee from the hips. Rather than the intact side simply
compensating for the impaired side in a diffuse manner, the operator
structure suggests that stroke channels transient coupling through a
knee-centered pathway on the non-paretic limb. That picture remains
consistent with the clinical observation that the ``good'' leg works
harder during stroke gait, but it localises the asymmetry to a specific
operator-level route through which perturbations are absorbed and
redistributed.

\textbf{Relationship to existing stability measures.} Bruijn et al.
(2013) evaluated gait stability measures against four validity criteria:
construct validity, convergent validity, predictive validity, and
discriminant validity. Floquet multipliers showed poor construct
validity for human gait--not only because strict periodicity is
violated, but because they require prohibitively large datasets
(\textgreater10,000 strides in Revzen and Guckenheimer (2008)
vs.~\textasciitilde15--50 strides in typical human trials), demand a
reliable phase estimator (Revzen and Guckenheimer 2008, 2012), and are
acutely noise-sensitive. \(\lambda_S\) showed good predictive validity
but limited discriminant validity between populations, and as Dingwell
and Kang (2007) demonstrated, the coexistence of positive \(\lambda_S\)
(local divergence) with orbital stability reflects the driven nature of
the system rather than a genuine autonomous instability. The
operator-level metrics introduced here show substantial discriminant
validity, with the strongest and most consistent separations in
AB-vs-stroke contrasts, and capture facets of locomotor
organisation--transient richness, frequency bandwidth, directional
coupling--that no single scalar can represent. Like DMD itself, these
metrics inherit the autonomous-model assumption, but they make the
strain of that assumption visible: non-normality is high precisely where
the linear autonomous model must work hardest to approximate driven,
nonlinear dynamics. Establishing their predictive validity for fall risk
and their convergent validity with \(\lambda_S\) are natural directions
for future work.

\textbf{Methodological considerations.} Our approach builds on the DMD
framework of Schmid (2010) and Tu et al. (2014), with delay embedding
following Takens (1981). Wu, Brunton, and Revzen (2021) showed that DMD
applied to limit-cycle data faces challenges from Krylov space closure
and non-normality-amplified estimation errors, and characterised
performance via a dataset metric (number of initial conditions
\(\times\) trajectory length). Our dataset of 299 trials with 1500 time
points each provides substantial coverage; the stability enforcement
(eigenvalue projection inside the unit disk) and SVD truncation
(retaining \textgreater99\% variance) guard against the most severe
numerical pathologies they identify. Nevertheless, several limitations
apply. The analysis uses a single treadmill dataset; overground walking
may produce different operator structure. The choice of embedding lag
(\(\tau = 10\)) is motivated by stride-period considerations at 100 Hz
but has not been exhaustively optimised. The linear operator is a
surrogate--it captures the Koopman-linearised dynamics faithfully only
within the delay-embedded subspace, and nonlinear effects beyond this
approximation are not modelled.

\textbf{Clinical implications.} Collectively, these results suggest that
post-stroke locomotor impairment is not a scalar degradation but a
structured reorganisation of transient dynamics. Restoring walking speed
or step symmetry may not restore the non-normal richness that
characterises healthy gait. The group-specific signatures identified
here--fragile non-normality, narrowed bandwidth in HF, and knee-centered
asymmetry that is strongest in LF--point toward different rehabilitation
targets for different functional levels. Interventions that broaden the
transient repertoire, such as perturbation-based balance training,
multi-frequency rhythmic cueing, or assistive devices that modulate
inter-limb coupling directionality, may prove more effective than
approaches targeting kinematic outcomes alone.

\textbf{Outlook.} The present work demonstrates that non-normality in
data-driven linear operators is not merely a numerical nuisance but a
window into the transient organisation of biological movement. Future
extensions--incorporating residual DMD for rigorous spectral convergence
guarantees (Colbrook, Ayton, and Szöke 2023), measure-theoretic Koopman
formulations (Mezić 2005, 2013), or hybrid embeddings that combine delay
coordinates with nonlinear observables--promise richer operator-level
phenotypes for clinical stratification and closed-loop control.

\textsubscript{Source:
\href{https://Michaelhess17.github.io/Synology/index.qmd.html}{Article
Notebook}}

\protect\phantomsection\label{refs}
\begin{CSLReferences}{1}{0}
\bibitem[\citeproctext]{ref-bruijn2013}
Bruijn, Sjoerd M., Onno G. Meijer, Peter J. Beek, and Jaap H. van Dieën.
2013. {``Assessing the Stability of Human Locomotion: A Review of
Current Measures.''} \emph{Journal of The Royal Society Interface} 10
(83): 20120999. \url{https://doi.org/10.1098/rsif.2012.0999}.

\bibitem[\citeproctext]{ref-brunton2017}
Brunton, Steven L., and J. Nathan Kutz. 2019. \emph{Data-Driven Science
and Engineering: Machine Learning, Dynamical Systems, and Control}. 1st
ed. Cambridge University Press.
\url{https://doi.org/10.1017/9781108380690}.

\bibitem[\citeproctext]{ref-colbrook2021}
Colbrook, Matthew J., Lorna J. Ayton, and Máté Szöke. 2023. {``Residual
{Dynamic} {Mode} {Decomposition}: Robust and Verified {Koopmanism}.''}
\emph{Journal of Fluid Mechanics} 955: A21.
\url{https://doi.org/10.1017/jfm.2022.1052}.

\bibitem[\citeproctext]{ref-dingwell2000}
Dingwell, Jonathan B., and Joseph P. Cusumano. 2000. {``Nonlinear Time
Series Analysis of Normal and Pathological Human Walking.''}
\emph{Chaos} 10 (4): 848--63. \url{https://doi.org/10.1063/1.1324008}.

\bibitem[\citeproctext]{ref-dingwell2006}
Dingwell, Jonathan B., and Hyun Gu Kang. 2007. {``Differences Between
Local and Orbital Dynamic Stability of Human Walking.''} \emph{Journal
of Biomechanical Engineering} 129 (4): 586--93.
\url{https://doi.org/10.1115/1.2746383}.

\bibitem[\citeproctext]{ref-henrici1962}
Henrici, Peter. 1962. {``Bounds for Iterates, Inverses, Spectral
Variation and Fields of Values of Non-Normal Matrices.''}
\emph{Numerische Mathematik} 4: 24--40.
\url{https://doi.org/10.1007/BF01386294}.

\bibitem[\citeproctext]{ref-mezic2005}
Mezić, Igor. 2005. {``Spectral Properties of Dynamical Systems, Model
Reduction and Decompositions.''} \emph{Nonlinear Dynamics} 41: 309--25.
\url{https://doi.org/10.1007/s11071-005-2824-x}.

\bibitem[\citeproctext]{ref-mezic2013}
---------. 2013. {``Analysis of Fluid Flows via Spectral Properties of
the {Koopman} Operator.''} \emph{Annual Review of Fluid Mechanics} 45:
357--78. \url{https://doi.org/10.1146/annurev-fluid-011212-140652}.

\bibitem[\citeproctext]{ref-revzen2008}
Revzen, Shai, and John M. Guckenheimer. 2008. {``Estimating the Phase of
Synchronized Oscillators.''} \emph{Physical Review E} 78 (5): 051907.
\url{https://doi.org/10.1103/PhysRevE.78.051907}.

\bibitem[\citeproctext]{ref-revzen2012}
---------. 2012. {``Finding the Dimension of Slow Dynamics in a Rhythmic
System.''} \emph{Journal of The Royal Society Interface} 9 (70):
957--71. \url{https://doi.org/10.1098/rsif.2011.0431}.

\bibitem[\citeproctext]{ref-schmid2010}
Schmid, Peter J. 2010. {``Dynamic Mode Decomposition of Numerical and
Experimental Data.''} \emph{Journal of Fluid Mechanics} 656: 5--28.
\url{https://doi.org/10.1017/S0022112010001217}.

\bibitem[\citeproctext]{ref-takens1981}
Takens, Floris. 1981. {``Detecting Strange Attractors in Turbulence.''}
In \emph{Dynamical Systems and Turbulence, Warwick 1980}, 898:366--81.
Lecture Notes in Mathematics. Springer.
\url{https://doi.org/10.1007/BFb0091924}.

\bibitem[\citeproctext]{ref-ting2007}
Ting, Lena H., and J. Lucas McKay. 2007. {``Neuromechanics of Muscle
Synergies for Posture and Movement.''} \emph{Current Opinion in
Neurobiology} 17 (6): 622--28.
\url{https://doi.org/10.1016/j.conb.2008.01.002}.

\bibitem[\citeproctext]{ref-trefethen2005}
Trefethen, Lloyd N., and Mark Embree. 2005. \emph{Spectra and
Pseudospectra: The Behavior of Nonnormal Matrices and Operators}.
Princeton University Press.

\bibitem[\citeproctext]{ref-tu2014}
Tu, Jonathan H., Clarence W. Rowley, Dirk M. Luchtenburg, Steven L.
Brunton, and J. Nathan Kutz. 2014. {``On {Dynamic} {Mode}
{Decomposition}: {Theory} and {Applications}.''} \emph{Journal of
Computational Dynamics} 1 (2): 391--421.
\url{https://doi.org/10.3934/jcd.2014.1.391}.

\bibitem[\citeproctext]{ref-winner2023}
Winner, Tristan, Michael C. Rosenberg, Trisha M. Kesar, Lena H. Ting,
and Gordon J. Berman. 2023. {``Discovering Individual-Specific Gait
Signatures from Data.''} \emph{bioRxiv}.
\url{https://doi.org/10.1101/2023.01.11.523590}.

\bibitem[\citeproctext]{ref-wu2021}
Wu, Jiahao, Steven L. Brunton, and Shai Revzen. 2021. {``Challenges in
{Dynamic} {Mode} {Decomposition}.''} \emph{Journal of The Royal Society
Interface} 18 (185): 20210686.
\url{https://doi.org/10.1098/rsif.2021.0686}.

\end{CSLReferences}




\end{document}
